\documentclass[10pt,a4paper]{article}
% ===================== PAQUETES =====================
\usepackage[utf8]{inputenc}
\usepackage[T1]{fontenc}
\usepackage[spanish,es-tabla]{babel}
\usepackage{amsmath,amssymb}
\usepackage{geometry}
\usepackage{enumitem}
\usepackage{hyperref}
\usepackage{fancyhdr}
\usepackage{titlesec}
\usepackage{setspace}
\usepackage{graphicx}
\usepackage{caption}
\usepackage{xcolor}
\usepackage{float}
\usepackage{pgfplots}
\pgfplotsset{compat=1.18}
\usepackage{pdflscape}
\usepackage{tabularx}
\usepackage{array}
\usepackage{booktabs}

% ===================== CONFIGURACIÓN =====================
\geometry{margin=2.5cm}
\setstretch{1.15}
\hypersetup{
    colorlinks=true,
    linkcolor=black,
    urlcolor=blue
}

\pagestyle{fancy}
\fancyhf{}
\captionsetup{font=small}
\setlength{\headheight}{32pt}

% FECHA
\usepackage{datetime}
\newdateformat{monthyeardate}{%
  \monthname[\THEMONTH], \THEYEAR}

% Header
\lhead{
    \includegraphics[width=3.6cm]{Logos/header_utfsm.png}
}
\rhead{
    \includegraphics[width=1.3cm]{Logos/Logo_industrias.png}
}

% Footer
\newcommand{\theCampus}{Vitacura}
\lfoot{\footnotesize ICN-292 — Sistemas de Información para la Gestión}
\cfoot{\footnotesize \thepage}

% Líneas
\renewcommand{\headrulewidth}{0.4pt}
\renewcommand{\footrulewidth}{0.4pt}

\usepackage{setspace}

\setlength{\parindent}{0pt}
\setlength{\parskip}{2pt}
\usepackage{listings}
\setlength{\headheight}{40pt}

\lstdefinestyle{sqlLaboratorio}{
    basicstyle=\ttfamily\footnotesize,
    columns=fullflexible,
    keepspaces=true,
    showstringspaces=false,
    breaklines=true,
    frame=single,
    framerule=0.35pt,
    rulecolor=\color{black!40},
    backgroundcolor=\color{black!2},
    aboveskip=5pt,
    belowskip=7pt,
    xleftmargin=3pt,
    xrightmargin=3pt,
    literate={á}{{\'a}}1 {é}{{\'e}}1 {í}{{\'i}}1 {ó}{{\'o}}1 {ú}{{\'u}}1 {ñ}{{\~n}}1 {Á}{{\'A}}1 {É}{{\'E}}1 {Í}{{\'I}}1 {Ó}{{\'O}}1 {Ú}{{\'U}}1
}
\lstset{style=sqlLaboratorio}

\begin{document}
\begin{center}
\vspace{-2mm}
{\Huge Laboratorio 4 --- SQL (1/2)}\\
{\Large \sc{ICN-292 --- Sistemas de Información para la Gestión}}
\vspace{2mm}
\hrule
\textbf{Charles Piket Silva}\\
{\small Ingeniería Civil Industrial, 202360512-3, \emph{cpiket@usm.cl}}\\
{\small \textbf{Repositorio GitHub:} \texttt{https://github.com/Cpikets/ICN292-Lab4-Piket-Charles.git}}
\end{center}
\vspace{5mm}

\section*{Consulta 1}
\begin{lstlisting}
SELECT nombre, genero
FROM autores
WHERE pais <> 'Chile' OR pais IS NULL
ORDER BY nombre;
\end{lstlisting}

La comparación \texttt{pais <> 'Chile'} no incluye los registros con \texttt{NULL}, porque su resultado es desconocido. Por eso se agrega \texttt{OR pais IS NULL}.

\begin{figure}[H]
    \centering
    \includegraphics[width=0.97\linewidth,keepaspectratio]{Evidencia_lab4/Consulta_01.png}
    \caption{Autores cuyo país no es Chile, incluyendo al autor sin país registrado.}
\end{figure}

\clearpage
\section*{Consulta 2}
\begin{lstlisting}
SELECT titulo, paginas
FROM libros
WHERE paginas BETWEEN 190 AND 220
ORDER BY paginas, titulo;
\end{lstlisting}

\begin{figure}[H]
    \centering
    \includegraphics[width=0.97\linewidth,keepaspectratio]{Evidencia_lab4/Consulta_02.png}
    \caption{Libros con 190 a 220 páginas, ordenados de menor a mayor.}
\end{figure}


\section*{Consulta 3}
\begin{lstlisting}
SELECT COUNT(*) AS cantidad_libros,
       SUM(ejemplares_vendidos) AS total_ejemplares_vendidos,
       ROUND(AVG(paginas), 2) AS promedio_paginas,
       MIN(anio) AS anio_mas_antiguo,
       MAX(ejemplares_vendidos) AS mayor_venta_individual
FROM libros;
\end{lstlisting}

\begin{figure}[H]
    \centering
    \includegraphics[width=0.97\linewidth,keepaspectratio]{Evidencia_lab4/Consulta_03.png}
    \caption{Indicadores de los doce libros de la carga oficial.}
\end{figure}

\clearpage
\section*{Consulta 4}
\begin{lstlisting}
SELECT DISTINCT genero
FROM autores
ORDER BY genero;
\end{lstlisting}

\begin{figure}[H]
    \centering
    \includegraphics[width=0.97\linewidth,keepaspectratio]{Evidencia_lab4/Consulta_04.png}
    \caption{Géneros distintos de los autores, ordenados alfabéticamente.}
\end{figure}


\section*{Consulta 5}
\begin{lstlisting}
SELECT a.nombre, COUNT(l.id_libro) AS libros_2024
FROM autores a
JOIN libros l ON a.id_autor = l.id_autor
WHERE l.anio = 2024
GROUP BY a.id_autor, a.nombre
HAVING COUNT(l.id_libro) > 1
ORDER BY a.nombre;
\end{lstlisting}

\begin{figure}[H]
    \centering
    \includegraphics[width=0.97\linewidth,keepaspectratio]{Evidencia_lab4/Consulta_05.png}
    \caption{Autores que publicaron más de un libro en 2024.}
\end{figure}

\clearpage
\section*{Consulta 6}
\begin{lstlisting}
SELECT titulo, ejemplares_vendidos
FROM libros
WHERE anio = 2024 AND ejemplares_vendidos > 2000
ORDER BY ejemplares_vendidos DESC
LIMIT 3;
\end{lstlisting}

\begin{figure}[H]
    \centering
    \includegraphics[width=0.97\linewidth,keepaspectratio]{Evidencia_lab4/Consulta_06.png}
    \caption{Tres libros de 2024 con mayores ventas entre los que superan 2.000 ejemplares.}
\end{figure}

\clearpage
\section*{Consulta 7}
\begin{lstlisting}
INSERT INTO libros (id_autor, titulo, paginas, ejemplares_vendidos, anio)
VALUES (3, 'Papel Mojado', 260, 50, 2024);
\end{lstlisting}

\begin{lstlisting}
SELECT titulo, paginas, ejemplares_vendidos, anio
FROM libros
WHERE titulo = 'Papel Mojado';
\end{lstlisting}

\begin{lstlisting}
DELETE FROM libros
WHERE titulo = 'Papel Mojado';
\end{lstlisting}

\begin{lstlisting}
SELECT titulo, paginas, ejemplares_vendidos, anio
FROM libros
WHERE titulo = 'Papel Mojado';
\end{lstlisting}

\begin{figure}[H]
    \centering
    \includegraphics[width=0.97\linewidth,keepaspectratio]{Evidencia_lab4/Consulta_07_antes.png}
    \caption{Registro de Papel Mojado después de insertarlo y antes de eliminarlo.}
\end{figure}

\begin{figure}[H]
    \centering
    \includegraphics[width=0.97\linewidth,keepaspectratio]{Evidencia_lab4/Consulta_07_despues.png}
    \caption{Consulta posterior al DELETE: Papel Mojado ya no aparece en el resultado.}
\end{figure}

\clearpage
\section*{Consulta 8}
\begin{lstlisting}
SELECT titulo, ejemplares_vendidos
FROM libros
WHERE titulo = 'Cielo de Enero';
\end{lstlisting}

\begin{lstlisting}
UPDATE libros
SET ejemplares_vendidos = ejemplares_vendidos + 100
WHERE titulo = 'Cielo de Enero';
\end{lstlisting}

\begin{lstlisting}
SELECT titulo, ejemplares_vendidos
FROM libros
WHERE titulo = 'Cielo de Enero';
\end{lstlisting}

\begin{figure}[H]
    \centering
    \includegraphics[width=0.97\linewidth,keepaspectratio]{Evidencia_lab4/Consulta_08_antes.png}
    \caption{SELECT previo al UPDATE de Cielo de Enero.}
\end{figure}

% La captura Consulta_08_update.png se incorporara cuando se reciba la evidencia real.

\begin{figure}[H]
    \centering
    \includegraphics[width=0.97\linewidth,keepaspectratio]{Evidencia_lab4/Consulta_08_despues.png}
    \caption{SELECT posterior al UPDATE: Cielo de Enero registra 1.600 ejemplares.}
\end{figure}

\clearpage
\section*{Consulta 9}
\begin{lstlisting}
INSERT INTO colecciones (titulo, fecha_publicacion) VALUES
('Horizontes Australes', '2024-01-15'),
('Letras del Mundo', '2023-06-20');
\end{lstlisting}

\begin{lstlisting}
UPDATE libros
SET id_coleccion = 1
WHERE id_libro IN (1, 2, 3, 6, 7, 8, 9, 10);
\end{lstlisting}

\begin{lstlisting}
UPDATE libros
SET id_coleccion = 2
WHERE id_libro IN (4, 5, 11, 12);
\end{lstlisting}

\begin{lstlisting}
SELECT l.titulo AS titulo_libro,
       c.titulo AS titulo_coleccion,
       c.fecha_publicacion
FROM libros l
JOIN colecciones c ON l.id_coleccion = c.id_coleccion
ORDER BY l.titulo;
\end{lstlisting}

\begin{figure}[H]
    \centering
    \includegraphics[width=0.97\linewidth,keepaspectratio]{Evidencia_lab4/Consulta_09.png}
    \caption{Libros con sus colecciones y fechas de publicación.}
\end{figure}

\clearpage
\section*{Consulta 10}
\begin{lstlisting}
INSERT INTO pagos (id_autor, fecha_pago, monto) VALUES
(1, '2024-05-15', 1250000.00),
(1, '2024-11-10', 750000.00),
(3, '2024-08-25', 980000.00);
\end{lstlisting}

\begin{lstlisting}
SELECT a.nombre, p.fecha_pago, p.monto
FROM pagos p
JOIN autores a ON p.id_autor = a.id_autor
ORDER BY a.nombre, p.fecha_pago;
\end{lstlisting}

\begin{lstlisting}
SELECT a.nombre, SUM(p.monto) AS total_pagado
FROM pagos p
JOIN autores a ON p.id_autor = a.id_autor
GROUP BY a.id_autor, a.nombre
ORDER BY a.nombre;
\end{lstlisting}

% La captura Consulta_10_detalle.png se incorporara cuando se reciba la evidencia real.
\begin{figure}[H]
    \centering
    \includegraphics[width=0.97\linewidth,keepaspectratio]{Evidencia_lab4/Consulta_10_detalle.png}
    \caption{Pagos individuales a los autores, con fechas e importes..}
\end{figure}

\begin{figure}[H]
    \centering
    \includegraphics[width=0.97\linewidth,keepaspectratio]{Evidencia_lab4/Consulta_10_totales.png}
    \caption{Total pagado por autor.}
\end{figure}


\clearpage
\section*{Entregable 1 --- Creación y carga de la base}
\lstset{basicstyle=\ttfamily\scriptsize}
\begin{lstlisting}
PRAGMA foreign_keys = ON;

DROP TABLE IF EXISTS libros;
DROP TABLE IF EXISTS pagos;
DROP TABLE IF EXISTS colecciones;
DROP TABLE IF EXISTS autores;

CREATE TABLE autores (
    id_autor INTEGER PRIMARY KEY AUTOINCREMENT,
    nombre TEXT NOT NULL,
    pais TEXT,
    genero TEXT NOT NULL
           CHECK (genero IN ('Cómic', 'Ensayo', 'Infantil', 'Novela', 'Poesía'))
);

CREATE TABLE colecciones (
    id_coleccion INTEGER PRIMARY KEY AUTOINCREMENT,
    titulo TEXT NOT NULL,
    fecha_publicacion DATE NOT NULL
);

CREATE TABLE libros (
    id_libro INTEGER PRIMARY KEY AUTOINCREMENT,
    id_autor INTEGER NOT NULL,
    id_coleccion INTEGER,
    titulo TEXT NOT NULL,
    paginas INTEGER,
    ejemplares_vendidos INTEGER,
    anio INTEGER,
    FOREIGN KEY (id_autor) REFERENCES autores(id_autor),
    FOREIGN KEY (id_coleccion) REFERENCES colecciones(id_coleccion)
);

CREATE TABLE pagos (
    id_pago INTEGER PRIMARY KEY AUTOINCREMENT,
    id_autor INTEGER NOT NULL,
    fecha_pago DATE NOT NULL,
    monto DECIMAL(10,2) NOT NULL,
    FOREIGN KEY (id_autor) REFERENCES autores(id_autor)
);

INSERT INTO autores (nombre, pais, genero) VALUES
('Elena Pizarro', 'Chile', 'Novela'),
('Harold Finch', 'Inglaterra', 'Ensayo'),
('Rodrigo Bahamonde', 'Chile', 'Poesía'),
('Valentina Correa', 'Argentina', 'Infantil'),
('Nora Keller', 'Estados Unidos', 'Novela'),
('Kenji Arakawa', NULL, 'Cómic');

INSERT INTO libros (id_autor, titulo, paginas, ejemplares_vendidos, anio) VALUES
(1, 'Cielo de Enero', 210, 1500, 2024),
(1, 'Marea Baja', 185, 800, 2023),
(1, 'El Norte', 200, 2200, 2024),
(2, 'Estrellas que Caen', 240, 3000, 2022),
(2, 'Máquinas', 215, 450, 2024),
(3, 'Luz Artificial', 300, 1200, 2024),
(3, 'Ruido Blanco', 280, 600, 2023),
(4, 'El Barrio de Lila', 190, 5000, 2024),
(4, 'Antes del Amanecer', 205, 2700, 2024),
(4, 'La Calle 8', 175, 900, 2023),
(5, 'Flor de Hierro', 195, 4100, 2024),
(5, 'Hora Dorada', 220, 1800, 2022);
\end{lstlisting}

\begin{figure}[H]
    \centering
    \includegraphics[width=0.85\textwidth,keepaspectratio]{Evidencia_lab4/Modelo_ER_Austral.png}
    \caption{Modelo entidad--relación de la base de datos Austral Editores, con sus cuatro tablas y respectivas claves primarias y foráneas.}
    \label{fig:modelo_er_austral}
\end{figure}
\end{document}
